\subsection{An append-only index at fixed affine rank}
\label{sec:dynamicfixedrank}

This subsection proves the fixed-rank index of
Theorem~\ref{thm:main-fixed-rank}.  A grid gives the index when the
temperature and coordinate bounds are polynomial in the width.
The grid controls variation of scores inside a cell.  It need not
recover their complete order.  A new key changes one cell, while
a query combines the cell envelopes into one exact rejection
sampler.  Keeping an earlier coordinate chart when the affine
span grows avoids any migration of stored records.

\begin{theorem}[Dynamic fixed-rank attention]
\label{thm:dynamicfixedrank}
Fix $r\ge1$ and positive dyadic coordinate bounds $K,Q$.
Let $\beta\ge0$ be dyadic, or $\beta=\sqrt n$, and choose a
fixed dyadic bound $\beta\le\bar\beta\le2\beta$ when
$\beta>0$, and set $\bar\beta=0$ when $\beta=0$.
A stream appends dyadic key--value records
$(k_j,v_j)$ with $k_j\in[-K,K]^n$ and $v_j\in[-1,1]$.
Every prefix of keys has affine dimension at most $r$.
Queries $q\in[-Q,Q]^n$ are explicitly supplied dyadic vectors
and are made only at nonempty prefixes $T\ge1$.
All scalar encodings and fixed parameters have at most $b$ bits,
with dyadics represented by binary integer numerators and
denominators.
Put $B=16+b+\lceil\log_2(T+n+r+2)\rceil$ at cache length $T$.

There is a uniform exact attention-index sampler with
\[
 \begin{split}
 \text{worst-case append work}
     &=O_r((n+\log(T+2))B^3),\\
 \text{stored bits}&=O_r((T+1)nB),\\
 \text{expected query work}
     &=O_r((n+M_T+\log(T+2))B^4),
 \end{split}
\]
where
\[
 M_T\le\min\left\{T,
        (r+1)(5+32r\bar\beta QK)^r\right\}.
\]
The expected number of proposed members is at most $17/8$.
The output index has probabilities proportional to
$\exp(\beta\langle q,k_j\rangle/n)$.  A fresh value sign then
has the exact attention-coordinate mean.  The claim holds
conditionally on every adaptive history of legal records and
queries.  At $\beta=0$, uniform index sampling suffices.

For fixed $r,K,Q$ and $\beta=\sqrt n$, expected query work is
\[
 O_{r,K,Q}\bigl((n+n^{r/2}+\log(T+2))B^4\bigr).
\]
The construction never scans earlier records after an append,
including an append that enlarges the span of a large degenerate
prefill.
\end{theorem}

The theorem charges operations on the supplied finite records.
Section~\ref{sec:momentdecoder} also charges the construction of
neural keys and values.

\begin{lemma}[A bounded rational coordinate chart]
\label{lem:dynamicchart}
Let $U$ be an $n$ by $d$ dyadic matrix of full column rank,
$1\le d\le r$, with entries of encoding length $O(b)$.
In $O_r(nB^3)$ bit operations one can find $d$ distinct row
indices $I$ such that
\[
 A=U(U_I)^{-1},\qquad A_I=I_d,\qquad |A_{ij}|\le2.
\]
The coefficients have $O_r(B)$ bits and a common determinant
denominator.
\end{lemma}

\paragraph{Proof idea.}
A large reconstruction coefficient identifies a row replacement
that more than doubles the determinant.  Integer determinant
bounds limit the number of such replacements.  This is the
classical maximum-volume argument; the coefficient and
row-replacement identities appear, for example, in
Mikhalev and Oseledets~\cite{mikhalev2018maxvol}, Section~3.

\begin{proof}
Find any nonsingular row minor by Gaussian elimination.  Clear
the common dyadic denominator of $U$, obtaining an integer
matrix $E$.  Suppose $C=E(E_I)^{-1}$ has an entry $|C_{ij}|>2$.
Replace selected row $j$ by row $i$.  This is a distinct row,
because the coefficients of selected rows form the identity.
Multilinearity of the determinant gives
\[
 |\det(E_{I,\mathrm{new}})|
       =|C_{ij}|\,|\det(E_I)|>2|\det(E_I)|.
\]
Every nonzero integer determinant has magnitude at least one.
If entries of $E$ have magnitude at most $H_E$, every minor is
at most $d!H_E^d$.  Its logarithm is $O_r(B)$.  Therefore only
$O_r(B)$ replacements occur.  At each stage Cramer's rule or
fixed-dimensional elimination computes all $nd$ coefficients
in $O_r(nB^2)$ bits.  Every intermediate determinant and
coefficient has $O_r(B)$ bits.  When no offending entry remains,
return the current indices.  The adjugate formula gives the
common denominator, and multiplication by the selected minor
gives $A_I=I_d$.
\end{proof}

\begin{lemma}[A finite cell envelope]
\label{lem:dynamiccellenvelope}
Partition $T$ members into nonempty cells.  Suppose cell $C$ has
a proxy score $a_C$ and each actual score satisfies
$|a_j-a_C|\le1/4$ for $j\in C$.  Suppose proxy comparisons
are exact and shifted negative exponentials admit the finite
enclosures of Lemma~\ref{lem:newnegativeexp}.  Then all cell
weights can be prepared using one exponential enclosure per
cell, and exact sampling needs at most $17/8$ proposed members
in expectation.  Final acceptance comparisons have a geometric
precision tail.
\end{lemma}

\paragraph{Proof idea.}
One score represents a cell to constant multiplicative accuracy.
A dyadic floor prevents a very small envelope weight from
requiring unbounded initial precision in the acceptance coin.
Adding an approximation allowance to construct a rejection
envelope, then refining only ambiguous comparisons, is an
established exact-sampling technique; see Brassard, Devroye,
and Gravel~\cite{brassard2019remote}, Sections~2 and~4.

\begin{proof}
Let $a_0=\max_C a_C$, and write
\[
 w_j=e^{a_j-a_0-1/4},\quad Z=\sum_jw_j,\qquad
 e_C=e^{a_C-a_0},\quad W_0=\sum_C|C|e_C.
\]
Every defining exponent is nonpositive.  For $j\in C$,
$w_j\le e_C\le e^{1/2}w_j<2w_j$, so $Z\le W_0<2Z$.
A member of a cell whose proxy equals $a_0$ has
$w_j\ge e^{-1/2}>1/2$, giving $Z>1/2$.

Set $p=\lceil\log_2(64T)\rceil$.  Enclose each $e_C$ in a
dyadic interval of width $2^{-p-2}$, round the upper endpoint
up to the grid $2^{-p}$, and add $2^{-p}$.  The resulting
dyadic $u_C$ satisfies
\[
 u_C\ge2^{-p},\qquad e_C\le u_C\le e_C+3\,2^{-p}.
\]
Thus the actual envelope obeys
\[
 Z\le W:=\sum_C|C|u_C
 <2Z+\frac3{64}
 \le\frac{67}{32}Z<\frac{17}{8}Z.
\]
Draw a cell proportionally to $|C|u_C$, then a uniform member
$j$ in it, and accept with probability $w_j/u_C$.  Its accepted
mass is
\[
 \frac{|C|u_C}{W}\frac1{|C|}\frac{w_j}{u_C}
 =\frac{w_j}{W}.
\]
The success probability is $Z/W$.  Independent trials therefore
return $w_j/Z$ and use $W/Z<17/8$ proposals in expectation.

The dyadic cell weights have a common denominator and $O(B)$
bits; integer cumulative weights sample a cell exactly.  Since
$u_C\ge2^{-p}$, an enclosure of width $2^{-p-t-3}$ for $w_j$
suffices at acceptance precision $t$.  Such an enclosure costs
$O((B+t)^4)$, so
Lemma~\ref{lem:newprecisionstop} decides the acceptance coin exactly
in expected $O(B^4)$ bit work per proposal, with a geometric
precision tail.  This stopping argument charges the
actual work of every reached precision stage.
\end{proof}

\paragraph{Proof idea for Theorem~\ref{thm:dynamicfixedrank}.}
Selected ambient coordinates describe every key in a bounded
chart.  Their grid controls the score variation for every legal
query.  When the span grows, a new chart receives only future
records.  There are at most $r+1$ charts, so querying all of them
has a bounded rank-dependent overhead.

\begin{proof}[Proof of Theorem~\ref{thm:dynamicfixedrank}]
Fix the first key $c$.  Maintain a basis $U$ of all differences
$k_j-c$ encountered so far.  An epoch starts at rank zero and
whenever this rank increases.  An epoch of dimension $d>0$
uses the chart of Lemma~\ref{lem:dynamicchart}, with
\[
 c_0=c-Ac_I,\qquad k=c_0+Ak_I.
\]
The selected coordinates $y=k_I$ lie in $[-K,K]^d$.
For a query let $v=q^TA$.  Since $|A_{ij}|\le2$,
$|v_i|\le2nQ$.  Choose
\[
 m=\left\lceil\log_2\max\{1,8r\bar\beta QK\}\right\rceil,
 \qquad \Delta=K2^{-m}.
\]
The integer cell tuple is $\lfloor y_i/\Delta\rfloor$.
Two members in one cell have coordinate differences less than
$\Delta$, so their scores differ by at most
\[
 \frac\beta n\sum_{i=1}^d|v_i|\Delta
 \le2\beta rQ\Delta\le\frac14.
\]
In one coordinate there are at most $2^{m+1}+1$ possible cell
indices.  This is at most $5+32r\bar\beta QK$.
The rank-zero epoch has one cell.

\paragraph{Appends.}
At an insertion test membership in the current affine span by
the exact equality $k-c=A(k_I-c_I)$.  This costs $O_r(nB^2)$.
If it holds, insert in the current epoch.  If it fails, append
$k-c$ to $U$, build its new bounded chart, and create a new
epoch containing only the new record.  Earlier epochs and their
records remain untouched.  At most $r$ increases are possible,
but even an individual increase costs only $O_r(nB^3)$.

Each epoch stores its occupied cells in a balanced search tree
keyed by their integer tuples.  A cell stores one representative
identifier and an appendable array of member identifiers.
For completeness, the latter can be a binary tree whose leaves
have local indices $0,\ldots,N-1$.  At a power-of-two capacity
boundary attach a new root and insert along the binary path
for the new local index.  There are $O(N)$ stored nodes and
$O(\log(N+2))$ pointer operations per append or random access.
All pointers, counters, and cell coordinates have $O_r(B)$
bits.  Drawing a uniform integer in $[N]$ by rejection from
binary bits takes expected $O(B^2)$ bit work.  These structures
give the stated worst-case append cost.

\paragraph{Queries.}
At a query compute $\langle q,c_0\rangle$ and $q^TA$ for each
epoch.  The number of epochs is at most $r+1$.  The entries of
$A$ have one common determinant denominator.  Combining that
denominator with the common dyadic query denominator shows that
an $n$-term projection adds only $O(\log n)$ numerator bits.
Each projection therefore has $O_r(B)$ encoding length and
costs $O_r(nB^2)$ bit work.  Evaluate each cell representative's
score from its $d$ selected coordinates, including the affine
offset so that scores in different epochs are comparable.

Apply Lemma~\ref{lem:dynamiccellenvelope} to the union of these
cells.  Their total number is at most
\[
 \min\left\{T,\sum_{d=0}^r(5+32r\bar\beta QK)^d\right\},
\]
which gives the displayed $M_T$.  At positive temperature,
comparisons of representative scores reduce to rational
comparisons before multiplication by $\beta/n$.  At standard
scaling, all shifted exponents are a rational plus a rational
multiple of $1/\sqrt n$, within
Lemma~\ref{lem:newnegativeexp}.  Thus one length-$n$ projection
per epoch, one enclosure per occupied cell, and a constant
expected number of member enclosures give the query bit bound.
Full records occupy $O(Tnb)$ bits, cell structures $O_r(TB)$,
and the bounded number of charts $O_r(nB)$.

\paragraph{Adaptive histories.}
Finally condition on any realized history of preceding legal
records and queries.  The current partition and its score
formulas are fixed.  Fresh random bits give exactly the law
proved by the envelope lemma.  Induction gives the claim for
adaptive histories.  At $\beta=\sqrt n$, a power-of-two upper
bound $\bar\beta<2\sqrt n$ gives $M_T=O_{r,K,Q}(n^{r/2})$.
\end{proof}

\begin{corollary}[Dynamic rank one at every temperature]
\label{cor:dynamicrankone}
For an append-only cache of affine rank at most one, numerical
temperature can be removed from the preceding complexity bound.
Worst-case append work is $O((n+\log(T+2))B^3)$ and expected
query work is $O((n+\log(T+2))B^4)$ at every allowed temperature.
The expected proposal count is at most $17/8$.
\end{corollary}

\paragraph{Proof idea.}
Maintain the scalar coordinate order with subtree member counts.
Rank selection replaces access to a stored permutation in the
static monotone envelope.

\begin{proof}
When the first nonconstant key arrives, choose a nonzero pivot
coordinate on its line.  Every preceding record has the same
key, so its existing member array becomes a single duplicate
group without scanning that array.  The rank promise prevents
any later change of line.  Maintain the ordered distinct pivot
coordinates in a balanced tree with subtree member counts.
Verify line membership and insert each new member exactly.

A query's one length-$n$ projection chooses the forward or
reverse order.  Retrieve the geometric-block endpoints of
Lemma~\ref{lem:newmonotoneenvelope} by rank selection and apply
the dyadic envelope of Theorem~\ref{thm:newrankone}.  One rank
selection costs $O(\log(T+2)B^2)=O(B^3)$ bits.  It is absorbed
by that endpoint's $O(B^4)$ exponential enclosure.  There are
$O(\log(T+2))$ endpoints and a constant expected number of
proposed ranks.  Hence the static query bound is retained,
including all exact arithmetic and precision refinement.
\end{proof}

\begin{corollary}[A certified neighborhood of a known key space]
\label{cor:dynamicapproximatespace}
Suppose a fixed rational chart $c_0+Ay$ is known,
$A_I=I_r$, $(c_0)_I=0$, and $|A_{ij}|\le2$.
The actual keys may have arbitrary rank, provided each satisfies
\[
 \|k-c_0-Ak_I\|_\infty\le\eta,
 \qquad \bar\beta Q\eta\le\frac18.
\]
The same append and query bounds hold with
$M_T\le\min\{T,(5+64r\bar\beta QK)^r\}$.
Only a constant expected number of proposed members require
length-$n$ actual-score evaluations.
\end{corollary}

\paragraph{Proof idea.}
Split the cell's score allowance between variation inside the
projected grid and the certified residual outside the key space.
The acceptance coin evaluates the actual score.

\begin{proof}
Use $m=\lceil\log_2\max\{1,16r\bar\beta QK\}\rceil$.
Projected scores in a cell vary by at most $1/8$.  The score
error from a residual of norm at most $\eta$ is at most
$\beta Q\eta\le1/8$.  Use the projected score of a stored
representative as the proxy.  Every actual member lies within
$1/4$ of it, so Lemma~\ref{lem:dynamiccellenvelope} applies.
Proxies use the fixed low-dimensional projection.  Actual
candidate scores use all $n$ coordinates, but there are fewer
than $17/8$ candidates in expectation.  The residual certificate
is checked by rational arithmetic at insertion.

If the promise is instead ordinary $\ell_\infty$ distance at
most $\eta_0$ from a known affine space, its bounded coordinate
chart has projection $P=AR_I$ with
$\|P\|_{\infty\to\infty}\le2r$.  Applying $I-P$ to an
approximating point gives residual at most $(1+2r)\eta_0$.
Thus $\bar\beta Q(1+2r)\eta_0\le1/8$ suffices.
\end{proof}

